\begin{afterword}
\summaryhead

The post continues the story of the previous day's conversation. The author put the Litany of Tarski to the woman and asked whether, if it were known that God does not exist, she should still believe in God. She hesitated and did not answer. The author suspects that God was not a part of her model of the world at all, only her belief in God. She then offered a counterexample to the litany: ``I believe that people are nicer than they really are.'' The author explained that saying something expresses belief in it and quoted Wittgenstein; she repeated the sentence, adding that she had put it that way for the author.

Later she said that nothing should be done differently if there were no God, and she paused for a long time when asked whether she was always surprised by people. The author concludes that her religion is the worship of her own failed attempt at self-deception. This also explains why the author argues that self-deception is hard: believing you have succeeded brings a placebo benefit, and the argument aims to remove it.

\responsehead

The post adds no new kind of evidence to the previous day's, and its conclusions grow firmer. ``So far as I can tell'' becomes ``I now realize that the whole essence of her philosophy was her belief that she had deceived herself,'' drawn from a pause, and then ``The attempt failed, but she is honestly unaware of this.'' The placebo benefit, offered the day before with ``I suppose,'' is now stated without a hedge. The evidence throughout is one conversation, as the author recalls it.

The post reads the exchange at its center one way only. Her sentence has the form of Moore's paradox. The author's gloss is that saying ``I believe people are nice'' means you believe you believe it. But the Wittgenstein line the author quotes comes from a section that takes a different view of the words ``I believe'': they are ``used like the assertion `This is the case'.'' She also said she had put the sentence that way ``for you.'' The post does not ask what she meant, although the wording it analyses was, by her account, chosen for her listener.

From this one case the post draws general claims. ``This is why intelligent people only have a certain amount of time'' to become atheists. A ``new sort of religion'' worships the worship of God, and the defenders of sanity must now fight it. Neither claim has support beyond the conversation.

The closing section states the author's aim. Arguing that self-deception is hard is described as ``taking direct aim at the placebo benefits'' of believing one has succeeded. Four days later the author went further (see the notes on ``Don't Believe You'll Self-Deceive'').

In short: firmer conclusions about one woman's mind from the same single conversation, a reading of her Moorean sentence that the quoted Wittgenstein does not share, and general claims about intelligent believers drawn from her case alone.
\end{afterword}
